\documentclass[11pt]{article}
\usepackage[margin=1in]{geometry}
\usepackage[T1]{fontenc}
\usepackage[utf8]{inputenc}
\usepackage{lmodern}
\usepackage{booktabs}
\usepackage{array}
\usepackage{longtable}
\usepackage{hyperref}
\usepackage{enumitem}
\setlist[itemize]{leftmargin=1.25em}
\setlength{\parskip}{0.5em}
\setlength{\parindent}{0pt}

\title{Milestone Report\\Multimodal Hallucination Alignment Project}
\author{Tanush Savadi}
\date{April 9, 2026}

\begin{document}
\maketitle

\section*{Project Goal}
The goal of this project is to study whether image aware preference optimization can reduce multimodal hallucinations better than standard DPO.

I am using Qwen2.5-VL-3B-Instruct as the base model.
I am training two aligned variants.
The first one is a standard DPO model.
The second one is an image aware DPO model.

The main idea is simple.
If the model is really grounded in the image then its answer should react in a reasonable way when I blank the image or swap the image.
That is what I am trying to test.

In simple words, this project is asking one main question.
If I fine tune a vision language model in a more image sensitive way, will it stop making up answers that are not supported by the picture.

\section*{Why This Project Matters}
Multimodal models can often sound confident even when the image does not support the answer.
That is a big problem if we want to trust them.
It is even worse when the answer sounds fluent and reasonable, because then it becomes harder to notice the mistake.

My proposal idea was to compare a standard preference learning method against a modified method that also checks whether the answer still looks good when the image is wrong.
So the project is not only about accuracy.
It is also about whether the model is truly using the image in a grounded way.

\section*{What The Project Actually Does}
The codebase has five main stages.

\begin{itemize}
    \item prepare the training and benchmark datasets into one clean format
    \item train a standard DPO model
    \item train an image aware DPO model
    \item evaluate both models on benchmark datasets
    \item save artifacts and later inspect them in a dashboard
\end{itemize}

The training data comes from a preference dataset.
Each example has an image, a prompt, a chosen answer, and a rejected answer.
So the model is learning which answer should be preferred.

The standard DPO model learns in the normal way.
It tries to prefer the chosen answer over the rejected answer.

The image aware DPO model does one extra thing.
It also looks at a mismatched image case.
So it is not just asking whether the chosen answer is better.
It is asking whether the chosen answer is better when the image actually matches the prompt.

For evaluation, I use three benchmark datasets.

\begin{itemize}
    \item HallusionBench for multimodal hallucination checks
    \item POPE for object existence yes and no questions
    \item ChartQA for chart based question answering
\end{itemize}

I also evaluate each sample under three image conditions.

\begin{itemize}
    \item original image
    \item blank image
    \item mismatched image
\end{itemize}

This lets me check not only if the answer is correct, but also if the answer depends on the image in a meaningful way.

\section*{What I Have Built So Far}
For this milestone I built the full pipeline.

\begin{itemize}
    \item data preparation for the training set and benchmark sets
    \item standard DPO training
    \item image aware DPO training
    \item benchmark evaluation on HallusionBench, POPE, and ChartQA
    \item artifact saving to Google Drive
    \item a Streamlit dashboard pipeline for later analysis
\end{itemize}

I also fixed a long list of engineering issues during this phase.
These were mostly multimodal batching bugs, Qwen image tensor issues, memory problems, and Colab runtime issues.
So one important milestone result is that the system now runs end to end at full pilot scale on remote GPU compute.

\section*{Experiment Setup}
The main full pilot config used:

\begin{itemize}
    \item train subset size: 8000
    \item validation size: 1000
    \item optimizer steps: 1000
    \item gradient accumulation steps: 8
    \item model: Qwen2.5-VL-3B-Instruct
    \item training methods: standard DPO and image aware DPO
    \item benchmarks: HallusionBench, POPE, and ChartQA
\end{itemize}

All long runs were saved to Google Drive so the artifacts stayed available even when the Colab session was unstable.

\section*{Progress So Far}
The project moved in stages.

First I got smoke runs working.
That gave me a small end to end sanity check for both the standard DPO and image aware DPO training paths.

After that I ran a reduced pilot during debugging.
That helped me fix practical issues without wasting a lot of Colab time.

After those stages worked, I switched back to the full pilot setup on stronger GPU hardware.
That full pilot is the main result of this milestone.

\section*{Compute Bottleneck}
One practical bottleneck in this project is runtime.
Even when using the Colab G4 GPU, the full pipeline is still slow.
Data preparation, full pilot training, and evaluation each take a long time.
So the main challenge is not only getting the code to work.
It is also making the experiments reliable enough to finish within Colab session limits.

This matters for the milestone because it affects how quickly I can test changes and rerun experiments.
It also means I have to be careful about saving artifacts to Google Drive and structuring the workflow so that partial progress is not lost if the session disconnects.

\section*{Completed Runs}
The following full pilot training runs finished successfully:

\begin{itemize}
    \item \texttt{2026-04-08-standard\_dpo-pilot-7}
    \item \texttt{2026-04-08-image\_aware\_dpo-pilot-7}
\end{itemize}

The following artifacts were saved for both runs:

\begin{itemize}
    \item adapter weights
    \item config file
    \item environment file
    \item training metrics
    \item preference preview parquet
\end{itemize}

For the standard DPO full pilot run, the evaluation artifacts were also written.
That means the standard baseline already has:

\begin{itemize}
    \item \texttt{predictions.jsonl}
    \item \texttt{dependence.jsonl}
    \item updated \texttt{metrics.json}
\end{itemize}

At the time of this milestone writeup the image aware evaluation was not fully finished yet.
So the final side by side comparison is still pending.

\section*{Training Results So Far}
The final training summaries from the two full pilot runs are shown below.

\begin{center}
\begin{tabular}{>{\raggedright}p{2.3cm} >{\raggedright}p{4.1cm} r r r r r r r}
\toprule
Model & Run ID & Global Step & Opt Steps & Last Loss & Last DPO & Matched Margin & Mismatched Margin & Gap Loss \\
\midrule
standard\_dpo & \texttt{2026-04-08-standard\_dpo-pilot-7} & 1000 & 1000 & 0.4772 & 0.4772 & 4.9171 & -- & -- \\
image\_aware\_dpo & \texttt{2026-04-08-image\_aware\_dpo-pilot-7} & 1000 & 1000 & 0.9165 & 0.7508 & -1.1211 & -0.7991 & 0.4220 \\
\bottomrule
\end{tabular}
\end{center}

I do not want to over claim from the raw training loss.
The two objectives are not the same.
So a lower standard DPO loss does not automatically mean it is the better model.
Still, the curves are useful for showing that both runs completed and trained in a stable way.

From the plots I made, the standard DPO run looks more numerically stable.
The image aware run is noisier, which makes sense because it includes the extra gap term.
The gap loss also shows that the image aware objective was active during training.

The training plots support that summary.
The standard DPO loss stayed in a tighter band and ended at a lower value.
The image aware run had a higher and noisier loss curve.
Its gap loss also had repeated spikes, which is expected because that model is actively trying to separate matched and mismatched image behavior.

So my current interpretation is not that the standard model is automatically better.
My interpretation is that the two models are learning different things.
The standard model looks easier to optimize.
The image aware model is doing a harder training job.

\section*{Preliminary Evaluation Results}
The standard DPO evaluation finished on all three benchmarks.
The most useful current numbers are below.

\begin{center}
\begin{tabular}{lll}
\toprule
Benchmark & Metric & Value \\
\midrule
ChartQA & relaxed accuracy & 0.3797 \\
ChartQA & augmented relaxed accuracy & 0.4990 \\
ChartQA & human relaxed accuracy & 0.2604 \\
HallusionBench & overall accuracy & 0.0177 \\
HallusionBench & illusion accuracy & 0.0417 \\
HallusionBench & video accuracy & 0.0824 \\
\bottomrule
\end{tabular}
\end{center}

These numbers show a mixed picture.
ChartQA is not great, but it is at least giving some useful signal.
HallusionBench is still very weak.
So right now the baseline model is clearly not solving multimodal hallucination well.

ChartQA is the clearest positive signal so far.
The overall relaxed accuracy is around 0.38.
The augmented split is better at around 0.50.
The human split is weaker at around 0.26.
So the model is getting some chart questions right, but the benchmark is still difficult.

HallusionBench is much worse.
The overall accuracy is only around 0.018.
Some subgroups are slightly better, like the video subgroup at about 0.082 and the illusion subgroup at about 0.042.
But overall the baseline is still struggling badly on this benchmark.

This is not a bad milestone result.
It actually tells me something important.
The problem is still hard for the baseline model.
That means there is real room for the image aware method to matter.

\section*{Dependence Analysis}
I also measured how often the model changed its answer when I changed the image input.
This was done using blank image and mismatched image variants.

\begin{center}
\begin{tabular}{lr}
\toprule
Metric & Value \\
\midrule
blank changed rate & 0.9667 \\
mismatch changed rate & 0.9285 \\
blank score drop mean & 0.0559 \\
mismatch score drop mean & 0.0522 \\
\bottomrule
\end{tabular}
\end{center}

This is one of the most interesting results so far.
The model changes its output a lot when the image is altered.
So it is clearly sensitive to the visual input.
That is good in one sense because it means the model is not totally ignoring the image.

But this does not automatically mean the model is well grounded.
A lot of the changed outputs are still wrong or only partly useful.
So my current view is that the model is image sensitive but still unreliable.

I think this is an important nuance.
If a model never changes when the image changes, then it is probably ignoring the image.
That is bad.
But if a model changes a lot and still gives bad answers, then it is using the image in an unstable or weak way.
That seems closer to what is happening here.

So the dependence analysis suggests visual sensitivity, but not strong visual reliability.

\section*{POPE Evaluation Issue}
During evaluation I found a bug in the POPE metric pipeline.
The model often answered in natural language like:

\begin{itemize}
    \item ``Yes, there is a snowboard in the image''
    \item ``There is no car present in the image''
\end{itemize}

The original POPE scorer only accepted exact \texttt{yes} or \texttt{no}.
Because of that, the first POPE results looked artificially bad.

I fixed this by normalizing common yes and no sentence patterns.
After that fix the standard DPO evaluation had to be rerun.
So for this milestone I am not treating the earlier POPE zero score as a real scientific result.

This was actually a useful milestone finding.
It showed that the evaluation pipeline needed more robust answer normalization.

This bug was important because the model was often saying sensible things like ``Yes, there is a snowboard in the image'' or ``There is no car present in the image''.
Those should count as yes and no answers.
So the original zero score was not a trustworthy measure of true model quality.

For the milestone, the safe thing is to report that the POPE evaluator needed a normalization fix and that the corrected rerun is part of the next step.

\section*{Qualitative Examples}
Looking at sample outputs was useful.
Some ChartQA examples were answered correctly with exact numeric outputs.
Other examples gave long descriptive answers when a short number was expected.
I also saw many cases where the answer changed under blank or mismatched images.

That supports the same high level story from the dependence summary.
The model does react to image changes.
But the reaction is not always reliable and it does not yet guarantee correct grounding.

Here are a few concrete examples from the standard DPO evaluation.

\begin{itemize}
    \item On a ChartQA sample asking ``Who portrayed Jon Snow'' the ground truth was ``Kit Harington'' but the model gave a long descriptive sentence and was marked wrong.
    \item On another ChartQA sample asking how many people applied for universal credit, the model answered exactly ``950000'' and was correct.
    \item On a wage question for Austria in 2010, the ground truth was ``42295'' but the model answered ``41542'' and was wrong.
\end{itemize}

These examples show two things.
First, the model can sometimes extract the right number exactly.
Second, it also often fails by giving a long fluent answer when a short exact answer is needed.

That is especially important for ChartQA.
The model is not only making factual mistakes.
Sometimes it is also failing the response format that the benchmark expects.

\section*{Main Takeaways For This Milestone}
The biggest thing I can say now is that the full project pipeline works.

\begin{itemize}
    \item full pilot training completed for both models
    \item full benchmark evaluation completed for the standard DPO baseline
    \item artifact saving and dashboard preparation work
    \item the baseline model still struggles on hard hallucination benchmarks
    \item the baseline model shows meaningful image sensitivity under dependence tests
    \item POPE scoring needed a normalization fix because the model answered in full sentences
    \item the image aware comparison is the next major result to finish
\end{itemize}

So this milestone is mainly about implementation and early experimental signal.
It is not the final comparison yet.
But it is enough to show that the project is real, working, and already producing useful evidence.

\section*{What I Will Do Next}
The next steps are:

\begin{itemize}
    \item finish the corrected standard DPO rerun for POPE sensitive reporting
    \item finish the image aware evaluation rerun
    \item compare standard DPO and image aware DPO side by side
    \item build the dashboard artifacts for both finished runs
    \item write up benchmark level and qualitative comparisons
    \item decide whether image aware DPO improves grounding or just changes output behavior
\end{itemize}

The final report will focus on that comparison.
Right now the baseline result already gives me a clear starting point.
The remaining task is to see whether the image aware method actually improves on it.

\section*{Honest Status At This Milestone}
I think this is a strong milestone because the full training pipeline is no longer hypothetical.
Both full pilot runs finished.
The baseline evaluation works.
The artifact system works.
The dependence analysis works.
And I already found one real evaluator bug through qualitative inspection of model outputs.

At the same time, I do not want to pretend the project is finished.
The most important final result is still the direct comparison between the standard DPO and image aware DPO models.
That comparison is close, but not fully closed out yet.

Another honest limitation is runtime.
Even with a stronger Colab GPU, the pipeline is still expensive in time.
That slows down iteration and makes debugging more costly than in a normal small ML class project.
So part of the engineering work in this milestone was not only model training.
It was also making the pipeline robust enough to survive long Colab runs and save progress correctly.

So my honest summary is this.
The implementation part is strong and complete enough to trust.
The experimental story is already meaningful.
The last step is to finish the corrected evaluation comparison and write the final conclusions from it.

\end{document}
